\documentclass[twocolumn]{aastex631}

\usepackage{amsmath}
\usepackage{multirow}
\usepackage{natbib}
\usepackage{graphicx} 
\usepackage{aas_macros}

\begin{document}

Lorentz-violating massive gravity theories present promising avenues for cosmological models but are often plagued by ghost instabilities and difficulties in reconciling with observational data. This paper introduces and explicitly constructs a novel Scalar-Modulated Internal Lorentz Symmetry (SMILS) within a specific class of such theories to address these critical issues. This composite, field-dependent symmetry, analytically derived by enforcing the invariance of the full Lagrangian in a Friedmann-Robertson-Walker background, involves coupled transformations of the physical metric, Stückelberg fields, and an auxiliary scalar field. Using a Hamiltonian analysis in the ADM formalism, we demonstrate that SMILS ensures the classical elimination of the Boulware-Deser ghost through the emergence of a unique second-class constraint system. Furthermore, a rigorous one-loop perturbative analysis, employing the background field method and examining the kinetic matrix eigenvalues of scalar perturbations, confirms that this symmetry robustly prevents the reappearance of ghost instabilities, a crucial aspect for theoretical viability. The SMILS imposes stringent constraints on the theory's potential functions, leading to a highly predictive subclass that intrinsically yields the speed of gravitational waves exactly equal to the speed of light, consistent with observations like GW170817. Moreover, by numerically solving the modified Friedmann equations and comparing with cosmic chronometer data, we show that the resulting theory successfully reproduces the observed cosmic expansion history. These findings establish a theoretically consistent and phenomenologically viable framework for Lorentz-violating massive gravity, robustly protected against ghosts from classical to one-loop levels and compatible with key cosmological observations.

\end{document}
                